\documentclass[master=ewtk]{kulemt}
\setup{
  title={Human-AI Co-Creation for Workplan Generation in Manufacturing: Optimizing Tool selection and Clamping strategies through Human-AI Dialogue},
  author={Jordaan Dimitrov},
  promotor={Prof.\,dr.\,ir.\ Dominiek Reynaerts \and Prof.\,dr.\,ir.\ Mathias Verbeke},
  assessor={Prof.\,dr.\,ir.\ Joost Duflou \and Dr.\,ir.\ Kerim Yildirim},
  assistant={Dr.\,ir.\ Ming Wu}}
% Remove the "%" on the next line for generating the cover page
%\setup{coverpageonly}
% Remove the "%" before the next "\setup" to generate only the first pages
% (e.g., if you are a Word user).
%\setup{frontpagesonly}

% If you want to include other LaTeX packages, do it here. 

% Finally the hyperref package is used for pdf files.
% This can be commented out for printed versions.
\usepackage[pdfusetitle,colorlinks,plainpages=false]{hyperref}
\usepackage[table]{xcolor} % 'table'optie is handig voor cell backgrounds
\usepackage{xltabular}    % combinatie van longtable + tabularx
\usepackage[table]{xcolor}
\usepackage{enumitem}
\usepackage{xltabular,enumitem} 
\usepackage[table]{xcolor} 
\usepackage{pdflscape}
%%%%%%%
% The lipsum package is used to generate random text.
% You never need this in a real master's thesis text!
\IfFileExists{lipsum.sty}%
 {\usepackage{lipsum}\SetLipsumDefault{11-13}}%
 {\newcommand{\lipsum}[1][11-13]{\par And some text: lipsum ##1.\par}}
%%%%%%%
\usepackage{longtable}
%%%%%%%
%\includeonly{chapter-n}
\begin{document}

\begin{preface}
This is my acknowledgment to thank everyone who has kept me busy.
I would like to thank my supervisor, my advisor, and the entire jury.
My family has also been very supportive, of course.
\end{preface}

\tableofcontents*

\begin{abstract}
  {\color{red} This \texttt{abstract} environment provides a summary of the work, which may or may not be extensive.
The intention is that this should be limited to 1 page.}

\begin{itemize}
	\item Autocam context, fixture planning
	\item human expertise dependence, brittle rule based systems
	\item aim is AI assisted clamping region detection
	\item method used is supervised learning, feedforward NN, face/mesh level processing
	\item training data obtained via manual labeling and parametric part generation
	\item processing done in multiple ways, cad face extraction, discretized mesh, adaptive mesh refinement
	\item surface/mesh information fed into NN: area; normals; outerness; opposing-surface validation; wall thickness 
	\item constraints are physical feasibility, reachability, total clamping force (thin walls)
	\item heavy computation via HPC cluster, training on gpu nodes, monitoring via 3 data sets (training, validation, xyz), one for tuning hyper parameters
	\item Metrics: precision; recall; F1-like combined score; false-positive suppression since that is more dangerous than missing a solution
	\item output gives user multiple solutions, can be ordered based on reachability, stability, etc
	\item challenges 
	\item future work
	
\end{itemize}


\end{abstract}

% A list of figures and tables is optional
%\listoffigures
%\listoftables
% If there are only a few figures and tables, it is better to use the following:
\listoffiguresandtables
% The list of symbols is also optional.
% This list must be created manually, e.g. as follows:
\chapter{List of abbreviations and symbols}
\section*{abbreviations}
\begin{flushleft}
  \renewcommand{\arraystretch}{1.1}
  \begin{tabularx}{\textwidth}{@{}p{12mm}X@{}}
    LoG   & Laplacian-of-Gaussian \\
    MSE   & Mean Square error \\
    PSNR  & Peak Signal-to-Noise ratio \\
  \end{tabularx}
\end{flushleft}
\section*{Symbols}
\begin{flushleft}
  \renewcommand{\arraystretch}{1.1}
  \begin{tabularx}{\textwidth}{@{}p{12mm}X@{}}
    42    & ``The Answer to the Ultimate Question of Life, the Universe,
            and Everything'' volgens de \cite{h2g2} \\
    $c$   & Speed of light \\
    $E$   & Energy \\
    $m$   & Mass \\
    $\pi$ & The number pi \\
  \end{tabularx}
\end{flushleft}

% Now comes the main text
\mainmatter

\chapter{Introduction}
\label{cha:intro}
The first contains a general introduction to the work. The goals are
defined and the modus operandi is explained.

\section{Overview AutoCAM project and thesis context}


\section{Motivation}

\section{Problem statement}

\section{Objectives}
\chapter{Background and related work}
\label{cha:1}
{\color{red} A chapter is a logical unit. It normally starts with an introduction, which
you are reading now. The last topic of the chapter holds the conclusion.}

\section{Current Fixture Planning Approaches}
First comes the introduction to this topic.

\subsection{Rule Based Generative fixture design}


\begin{itemize}
	\item yes no decision trees
	\item hardcoded geometric constraints
	\item deterministic logic means low flexibility
	\item brittle with more difficult geometry
	\item low adaptability
\end{itemize}

\begin{figure}
	\centering
	\includegraphics[scale=0.7]{images/Decision tree generative.png}
	\caption{Generative fixture design decision tree (source)}
	\label{fig:generative}
\end{figure}

\subsection{Case based variant fixture design}
\begin{itemize}
	\item Group technology classification (eg opitz)
	\item case library of prior part solutions
	\item similarity search
	\item adaptation step
	\item dataset quality determines performance
	\item limited innovation because biased toward historical design, unless adaptation step is very sophisticated
	\item adaptation hard for new features
	\item schematic \ref{fig:variant}
	 
	 
	 \begin{figure}
	 	\centering
	 	\includegraphics[scale=0.7]{images/variant.png}
	 	\caption{Variant fixture design scheme (source)}
	 	\label{fig:variant}
	 \end{figure}
	 
	 
	 
\end{itemize}

\section{AI in manufacturing}
\subsection{General role of AI in manufacturing}
\begin{itemize}
	\item data driven optimization
	\item reducing human dependency
\end{itemize}
\subsection{AI for CAD/CAM applications}
\begin{itemize}
	\item feature/pattern recognition
	\item automatic toolpath generation
	\item geometric classification (paper from presentation) \ref{fig:deeplearningclassification}
\end{itemize}

 \begin{figure}
	\centering
	\includegraphics[scale=0.4]{images/deeplearningclassification.png}
	\caption{Deep learning classification}
	\label{fig:deeplearningclassification}
\end{figure}


\subsection{Motivation for AI based methods}
\begin{itemize}
	\item abstraction beyond hard coded rules
	\item ability to handle many geometric inputs without breaking (robustness)
	\item probabilistic prediction
	\end{itemize}

\subsection{Challenges \& Limitations}
\begin{itemize}
	\item large labeled dataset needed
	\item computational cost
	\item explainibility issues (black box)
	\item safety concerns with decision making (liability?)
\end{itemize}
\subsection{Trends in Industry}
\begin{itemize}
	\item Move toward hybrid human-AI system, rather than fully autonomous decision making
	\item digital twins and AI monitoring
	\item predictive maintenance
\end{itemize}

\section{Conclusion}
\begin{itemize}
	\item classical methods rigid, brittle and have limited adaptability
	\item ai enables datadriven reasoning, robustness, abstraction, probablistic outputs
	\item but has its challenges such as data requirements, computational cost, explainability, safety
	\item industry trend: hybrid workflow
\end{itemize}
\chapter{Methodology}
\label{cha:2}

\section{Surface processing pipeline}

\begin{itemize}
	\item Cad model is input
	\item extraction surfaces 
	\item filtering out irrelevant geometry?
	\item computation of geometric descriptors (normals, curvature, area etc)
	\end{itemize}
	
	
	
	\begin{figure}
		\centering
		\includegraphics[scale=0.35]{images/NN visualization.jpg}
		\caption{Feedforward network visualization with two hidden layers}
		\label{fig:NNvisualization}
	\end{figure}
	
	
	
	
\begin{figure}
	\centering
	\includegraphics[scale=0.3]{images/pipeline faces.jpg}
	\caption{Pipeline for surface processing}
	\label{fig:surfacepipeline}
\end{figure}


\section{Discretized Geometry Processing}
\begin{figure}
	\centering
	\includegraphics[scale=0.3]{images/pipeline discretized.jpg}
	\caption{Discretized Pipeline}
	\label{fig:discretized pipeline}
\end{figure}

\begin{itemize}
	\item conversion of brep surfaces into discretized mesh
	\item adaptive mesh refinement
	\item computationally much heavier
	\item physical contraints into loss function: net clamping force zero, patch size constraint to prevent isloated micro triangles
	\item more uncertainty
\end{itemize}
\begin{figure}
	\centering
	\includegraphics[scale=0.7]{images/adaptivemeshrefinement.png}
	\caption{Adaptive mesh refinement}
	\label{fig:adaptive mesh }
\end{figure}

\section{Scaling model}

\begin{itemize}
	\item HPC cluster (reference them as they ask on their website)
	\item optimize model to run on GPU node (parallelization)
	\item parametric modeling to create a large number of parts
	\item used Blender with base part and modifications such as rotation scaling holes chamfers etc
\end{itemize}


\begin{figure}
	\centering
	\includegraphics[scale=0.5]{images/Parametric modelling.png}
	\caption{Parametrically generating parts}
	\label{fig:parametric modelling}
\end{figure}



\section{Feature enrichment}
\begin{itemize}
\item goal is to capture more geometric information to feed into model, therefore the following features were added
\end{itemize}

\subsection{Outerness}

\begin{figure}
	\centering
	\includegraphics[scale=0.7]{images/outerness.png}
	\caption{Outerness}
	\label{fig:outerness}
\end{figure}

\subsection{Opposing surface validation}

\begin{figure}
	\centering
	\includegraphics[scale=0.7]{images/opposing surface validation.jpg}
	\caption{Opposing surface validation}
	\label{fig:opposing surface validation}
\end{figure}


\subsection{Wall thickness}


\begin{figure}
	\centering
	\includegraphics[scale=0.7]{images/wall thickness.png}
	\caption{Wall Thickness}
	\label{fig:wall thickness}
\end{figure}


\section{Post processing}
\begin{itemize}
	\item AI model outputs a single colored region containing multiple solutions
	\item goal is to extract the different solutions and rank them based on feasibility, reachability, stability etcand provide a visualization
	\item The goal is not to automatically choose the best solution, it is a tool to help the process planner or operator by providing a suggestion which he/she should approve
\end{itemize}


\begin{figure}
	\centering
	\includegraphics[scale=0.3]{images/postprocessing.jpg}
	\caption{Visualization of different solutions on the same part in post processing}
	\label{fig:postprocessing}
\end{figure}


\section{Conclusion}


% ... and so on until
\chapter{Results}

\section{Surface processing}
\begin{itemize}
	\item Probability heat map
	\item Overfitting
	\item Labeling accuracy
\end{itemize}

\begin{figure}
	\centering
	\includegraphics[scale=0.8]{images/Overfitting example.png}
	\caption{Overfitting example}
	\label{fig:overfitting}
\end{figure}


\begin{figure}
	\centering
	\includegraphics[scale=0.3]{images/visualizationsurfaceprocessing.jpg}
	\caption{Visualization examples surface processing}
	\label{fig:visualizationsurfaceprocessing}
\end{figure}
\begin{figure}
	\centering
	\includegraphics[scale=0.5]{images/traininglosssurface.png}
	\caption{Visualization examples surface processing}
	\label{fig:trainglosssurfaceprocessing}
\end{figure}


\section{Discretized geometry processing}
\begin{figure}
	\centering
	\includegraphics[scale=0.7]{images/initial result discretized.png}
	\caption{First result discretized}
	\label{fig:initial result}
\end{figure}

\begin{figure}
	\centering
	\includegraphics[scale=0.7]{images/micro patch solution.png}
	\caption{Illustration of micropatch solution}
	\label{fig:micropatch solutions}
\end{figure}

\begin{figure}
	\centering
	\includegraphics[scale=0.9]{images/f1 score.png}
	\caption{F1 score result}
	\label{fig:f1 score result}
\end{figure}

\section{Conclusion}

\chapter{Discussion and Future Research Directions}
\label{cha:n}

\section{The First Topic of this Chapter}
\subsection{Item 1}
\subsubsection{Sub-item 1}

\subsubsection{Sub-item 2}

\subsection{Item 2}

\section{The Second Topic}

\section{Conclusion}

\include{conclusion}

% If you have appendices:
\appendixpage*          % if wanted
\appendix
\chapter{The First Appendix}
\label{app:A}
Appendices hold useful data which is not essential to understand the work
done in the master's thesis. An example is a (program) source.
An appendix can also have sections as well as figures and references\cite{h2g2}.

\section{Gant chart}

\begin{figure}
	\centering
	\includegraphics[scale=0.3\add{7}]{images/gant chart.jpg}
	\caption{Gant chart}
	\label{fig: gant chart}
\end{figure}
% ... and so on until
\include{appendix-n}
\backmatter
% The bibliography comes after the appendices.
% You can replace the standard "abbrv" bibliography style by another one.
\bibliographystyle{abbrv}
\bibliography{references}
\include{CoC}
\end{document}
